% !TEX program = lualatex
% Guia paso a paso para alumnos - DUOC FPY1101
% Compilar con: lualatex Guia-Alumno.tex (dos veces para TOC)
\DocumentMetadata{
  pdfstandard   = ua-2,
  pdfversion    = 2.0,
  lang          = es-CL,
  tagging       = on,
  testphase     = {phase-III,tagpdf,table,firstaid}
}

\documentclass[11pt,a4paper,oneside]{report}

\usepackage{fontspec}
\usepackage{polyglossia}
\setdefaultlanguage{spanish}

\usepackage[a4paper,margin=2.5cm,top=2.8cm,bottom=2.8cm]{geometry}
\usepackage{xcolor}

% Colores con contraste WCAG AA (>= 4.5:1 sobre fondo blanco)
\definecolor{primary}{HTML}{0050A0}     % azul oscuro
\definecolor{accent}{HTML}{8B4500}      % naranja oscuro
\definecolor{success}{HTML}{1B5E20}     % verde oscuro
\definecolor{danger}{HTML}{B71C1C}      % rojo oscuro
\definecolor{boxBg}{HTML}{F4F6FA}       % gris muy claro fondo cajas
\definecolor{boxBorder}{HTML}{C5D0E0}   % gris-azul borde
\definecolor{codeBg}{HTML}{F0F0F0}      % gris claro para codigo

\usepackage{hyperref}
\hypersetup{
  pdftitle    = {Guia del Alumno - Entrega de Evaluacion a GitHub},
  pdfauthor   = {DUOC UC - Roberto Arce - FPY1101},
  pdfsubject  = {Guia paso a paso para entrega de evaluaciones via GitHub Classroom},
  pdfkeywords = {GitHub, Classroom, evaluacion, DUOC, FPY1101, Python, IDLE},
  pdflang     = {es-CL},
  colorlinks  = true,
  linkcolor   = primary,
  urlcolor    = primary,
  citecolor   = primary,
  bookmarksnumbered = true
}

\usepackage{titlesec}
\titleformat{\chapter}[display]
  {\normalfont\Huge\bfseries\color{primary}}
  {\filright\Large\MakeUppercase{\chaptertitlename}\ \thechapter}
  {1ex}{\filright}
\titlespacing*{\chapter}{0pt}{-30pt}{20pt}

\titleformat{\section}
  {\normalfont\Large\bfseries\color{primary}}{\thesection}{1em}{}

\titleformat{\subsection}
  {\normalfont\large\bfseries\color{accent}}{\thesubsection}{1em}{}

\usepackage{enumitem}
\setlist{itemsep=2pt,parsep=2pt}

\usepackage{tcolorbox}
\tcbuselibrary{skins,breakable}

\newtcolorbox{notebox}[1][]{
  enhanced, breakable, colback=boxBg, colframe=boxBorder,
  boxrule=0.5pt, arc=3pt, left=8pt, right=8pt, top=6pt, bottom=6pt,
  title=#1, fonttitle=\bfseries\color{primary},
  coltitle=primary, colbacktitle=boxBg,
  attach boxed title to top left={xshift=8pt, yshift=-9pt},
  boxed title style={size=small, colback=white, colframe=white}
}

\newtcolorbox{warnbox}[1][]{
  enhanced, breakable, colback=boxBg, colframe=danger,
  boxrule=1pt, arc=3pt, left=8pt, right=8pt, top=6pt, bottom=6pt,
  title=#1, fonttitle=\bfseries\color{white},
  coltitle=white, colbacktitle=danger,
  attach boxed title to top left={xshift=8pt, yshift=-9pt},
  boxed title style={size=small, colback=danger, colframe=danger}
}

\newtcolorbox{successbox}[1][]{
  enhanced, breakable, colback=boxBg, colframe=success,
  boxrule=1pt, arc=3pt, left=8pt, right=8pt, top=6pt, bottom=6pt,
  title=#1, fonttitle=\bfseries\color{white},
  coltitle=white, colbacktitle=success,
  attach boxed title to top left={xshift=8pt, yshift=-9pt},
  boxed title style={size=small, colback=success, colframe=success}
}

\usepackage{listings}
\lstset{
  basicstyle=\ttfamily\small,
  backgroundcolor=\color{codeBg},
  frame=single,
  rulecolor=\color{boxBorder},
  framesep=4pt,
  xleftmargin=0pt,
  xrightmargin=0pt,
  breaklines=true,
  showstringspaces=false,
  literate={á}{{\'a}}1 {é}{{\'e}}1 {í}{{\'i}}1 {ó}{{\'o}}1 {ú}{{\'u}}1
           {ñ}{{\~n}}1 {Á}{{\'A}}1 {É}{{\'E}}1 {Í}{{\'I}}1
           {Ó}{{\'O}}1 {Ú}{{\'U}}1 {Ñ}{{\~N}}1
}

\usepackage{fancyhdr}
\pagestyle{fancy}
\fancyhf{}
\renewcommand{\headrulewidth}{0.4pt}
\renewcommand{\footrulewidth}{0pt}
\fancyhead[L]{\small\color{primary}Entrega de Evaluacion a GitHub}
\fancyhead[R]{\small\color{primary}Guia del Alumno}
\fancyfoot[C]{\thepage}

\renewcommand{\contentsname}{Indice}
\renewcommand{\chaptername}{Capitulo}

% ============================================================
\begin{document}
\sloppy

% ===== Portada =====
\begin{titlepage}
\centering
\vspace*{2cm}

{\color{primary}\Huge\bfseries Entrega de Evaluacion a GitHub}\\[8pt]
{\color{accent}\LARGE Guia paso a paso para alumnos}\\[2.5cm]

\rule{0.7\textwidth}{0.8pt}\\[10pt]
{\large DUOC UC\,---\,Asignatura FPY1101}\\[4pt]
{\large Programacion en Python}\\[10pt]
\rule{0.7\textwidth}{0.8pt}\\[3cm]

{\large Profesor: Roberto Arce}\\[8pt]
{\large Secciones 001D \textbullet{} 002D \textbullet{} 003D}\\[3cm]

{\small Version 1.2.x}\\[4pt]
{\small \href{https://github.com/Robbyfuu/entrega-evaluacion-github}{github.com/Robbyfuu/entrega-evaluacion-github}}\\

\vfill
\end{titlepage}

\tableofcontents
\clearpage

% ============================================================
\chapter{Resumen del proceso}

Esta guia te explica como entregar tus evaluaciones de programacion usando un script automatico que sube tu codigo a GitHub.

El proceso se divide en \textbf{dos momentos}:

\begin{notebox}[En tu casa\,---\,Antes del examen]
Una vez por semestre, debes preparar tu PC: instalar dependencias, crear tu cuenta de GitHub, iniciar sesion, elegir tu seccion y aceptar la tarea asignada por el profesor.

Luego, antes de cada evaluacion, repites solo el paso de \textbf{aceptar la tarea} de Classroom (un click).
\end{notebox}

\begin{successbox}[En el aula\,---\,Dia del examen]
Solo tres pasos: clonar el repositorio que ya aceptaste en casa, programar tu evaluacion en IDLE, y subir los archivos. El script entrega automaticamente el enlace que pegas en el AVA.
\end{successbox}

\section{Que vas a necesitar}

\begin{itemize}
  \item Windows 10 o 11
  \item Conexion a internet
  \item Una cuenta de GitHub (gratuita\,---\,la guia explica como crearla)
  \item El paquete \texttt{EntregaEvaluacion-Alumno.zip} que entrego el profesor
\end{itemize}

\section{Que NO necesitas saber}

\begin{itemize}
  \item No necesitas saber usar la consola
  \item No necesitas instalar git ni gh manualmente (el script lo hace solo)
  \item No necesitas memorizar comandos
\end{itemize}

% ============================================================
\chapter{En casa\,---\,Preparacion (una sola vez)}

Sigue estos pasos ANTES del dia de la evaluacion. Lo ideal es hacerlo el primer dia de clases para evitar problemas.

\section{Paso 1\,---\,Descargar y descomprimir el script}

\begin{enumerate}
  \item El profesor te entregara un archivo llamado \texttt{EntregaEvaluacion-Alumno.zip}.
  \item Descargalo y guardalo en una carpeta facil de encontrar (por ejemplo \texttt{Escritorio\textbackslash{}DUOC}).
  \item Click derecho sobre el ZIP $\rightarrow$ \textbf{Extraer todo}.
  \item Veras una carpeta con estos archivos:
\end{enumerate}

\begin{lstlisting}
EntregaEvaluacion-Alumno/
  Subir-Evaluacion.bat       <- Doble-click aqui
  Subir-Evaluacion.ps1
  Reset-GitHubAuth.ps1
  README.md
\end{lstlisting}

\section{Paso 2\,---\,Crear tu cuenta de GitHub (si no tienes)}

Si ya tienes cuenta, salta al Paso 3.

\begin{enumerate}
  \item Haz doble-click en \texttt{Subir-Evaluacion.bat}.
  \item Windows puede mostrar una advertencia de seguridad: click en \textbf{Mas informacion} $\rightarrow$ \textbf{Ejecutar de todas formas}.
  \item Se abre el formulario principal.
  \item En la esquina superior derecha veras el enlace: \emph{\textquotedblleft\,No tienes cuenta de GitHub? Creala aqui\,\textquotedblright}.
  \item Click en el enlace $\rightarrow$ se abre un dialogo con instrucciones $\rightarrow$ click \textbf{Abrir GitHub}.
  \item Se abre tu navegador en \url{https://github.com/signup}.
  \item Completa: email personal, contrasena (minimo 8 caracteres con letras y numeros), username.
  \item Verifica tu email con el codigo de 8 digitos que recibes.
  \item Listo. Vuelve al script y click \textbf{Ya tengo cuenta, iniciar sesion}.
\end{enumerate}

\begin{notebox}[Sugerencia de username]
Usa algo profesional como \texttt{tu-nombre-apellido} o \texttt{tu-nombre.duoc}. Lo veras toda tu carrera y aparece publicamente en GitHub.
\end{notebox}

\section{Paso 3\,---\,Instalar dependencias (winget)}

La primera vez que abres el script, te ofrecera instalar \texttt{git} y \texttt{gh} (GitHub CLI) usando \texttt{winget}.

\begin{enumerate}
  \item Cuando aparezca el dialogo \textbf{Instalar dependencias}, click \textbf{Si}.
  \item Acepta los permisos de administrador que pida Windows.
  \item Espera 1\,---\,2 minutos. Veras una consola negra instalando.
  \item Cuando termine, \textbf{cierra y vuelve a abrir} el script para que detecte las nuevas herramientas.
\end{enumerate}

\section{Paso 4\,---\,Seleccionar tu seccion}

La primera vez aparecera un dialogo \textbf{Selecciona tu seccion}.

\begin{itemize}
  \item Elige \textbf{001D}, \textbf{002D} o \textbf{003D} segun corresponda a tu inscripcion.
  \item Click \textbf{Continuar}.
\end{itemize}

\begin{warnbox}[Importante]
Si te equivocas de seccion, los assignments del Classroom no apareceran. Puedes corregirlo desde el dropdown \emph{\textquotedblleft\,Tu seccion\,\textquotedblright} arriba del formulario principal.
\end{warnbox}

\section{Paso 5\,---\,Iniciar sesion en GitHub}

\begin{enumerate}
  \item En el formulario principal, click en cualquier accion (por ejemplo, \textbf{Crear Repo}).
  \item Si no estas conectado, el script muestra un dialogo con un \textbf{codigo} grande tipo \texttt{XXXX-XXXX}.
  \item Abre en tu navegador (o celular): \url{https://github.com/login/device}
  \item Ingresa el codigo (boton \textbf{Copiar codigo} ayuda).
  \item Inicia sesion en GitHub si no estabas conectado.
  \item Autoriza el acceso a \emph{\textquotedblleft GitHub CLI\,\textquotedblright}.
  \item El script detecta automaticamente la autorizacion y continua.
\end{enumerate}

\begin{successbox}[Una sola vez]
Tu sesion queda guardada permanentemente en este PC. No tendras que volver a loguearte (a menos que cierres sesion manualmente).
\end{successbox}

\section{Paso 6\,---\,Aceptar la tarea de GitHub Classroom}

Esto es lo MAS importante para preparar el dia del examen:

\begin{enumerate}
  \item Cuando hay tareas activas, veras un banner azul arriba del formulario:\\
        \emph{\textquotedblleft 📚 Tienes 1 tarea(s) de Classroom\,---\,Click para aceptarlas\,\textquotedblright}
  \item Click en el banner.
  \item Se abre un dialogo con la lista de tareas (filtradas por tu seccion).
  \item Click en \textbf{Aceptar tarea} de la evaluacion correspondiente.
  \item Se abre tu navegador en GitHub Classroom.
  \item Verifica que la cuenta logueada en GitHub es la \textbf{tuya} (no la del PC del aula si estas en clases).
  \item Click \textbf{Accept this assignment} en GitHub.
  \item GitHub Classroom crea automaticamente un repositorio en tu cuenta con un nombre tipo \texttt{evaluacion-1-tu-username}.
  \item Cierra el navegador y vuelve al script.
\end{enumerate}

\begin{warnbox}[NO clones todavia]
Acepta la tarea pero NO la clones en casa. Si clonas y editas codigo en casa, el sistema lo detectara como TRAMPA cuando lo subas el dia del examen.

Solo acepta la tarea. El repositorio queda esperando en tu cuenta de GitHub, vacio.
\end{warnbox}

\section{Paso 7\,---\,Cerrar el script y guardar el PC}

Si trabajas desde tu PC personal, cierra el script normalmente y apaga. Todo queda configurado.

\textbf{El dia de la evaluacion solo tendras que abrir el script y hacer clone + edit + push.}

% ============================================================
\chapter{Dia del examen\,---\,En el aula}

Llegaste al aula. Estos son los pasos:

\section{Paso 1\,---\,Abrir el script}

\begin{enumerate}
  \item Si el PC del aula ya tiene el script: doble-click \texttt{Subir-Evaluacion.bat}.
  \item Si NO lo tiene: descargalo de nuevo del enlace que entrega el profesor y extrae el ZIP.
  \item En PCs del aula puede que tengas que repetir el login (Paso 5 del capitulo anterior).
\end{enumerate}

\section{Paso 2\,---\,Verificar tu seccion}

Mira el dropdown \emph{\textquotedblleft\,Tu seccion\,\textquotedblright} arriba: debe decir tu seccion correcta. Si no, cambiala.

\section{Paso 3\,---\,Clonar el repositorio}

\begin{enumerate}
  \item En el formulario principal, selecciona \textbf{Usar repositorio existente de mi cuenta}.
  \item Los campos Nombre y Tipo se deshabilitan. Aparece el dropdown \textbf{Tu repositorio}.
  \item Click en \textbf{Refrescar} para cargar la lista de tus repositorios.
  \item Selecciona el repositorio del assignment (tipo \texttt{[Pub] evaluacion-1-tu-username}).
  \item Click en \textbf{Clonar Repo}.
\end{enumerate}

El script:

\begin{itemize}
  \item Descarga tu repositorio a \texttt{Escritorio\textbackslash{}evaluacion-1-tu-username}
  \item Verifica que el repo este VACIO (anti-trampa)
  \item Abre \textbf{IDLE de Python} apuntando a esa carpeta
  \item Te muestra un mensaje con la ruta exacta
\end{itemize}

\section{Paso 4\,---\,Programar la evaluacion en IDLE}

\begin{enumerate}
  \item En IDLE, click \textbf{File} $\rightarrow$ \textbf{New File} para crear un archivo nuevo.
  \item Guarda como \texttt{evaluacion.py} (o el nombre que indique el profesor) DENTRO de la carpeta clonada.
  \item Programa tu solucion.
  \item Guarda con \texttt{Ctrl+S} cada vez que avances.
\end{enumerate}

\begin{notebox}[Probar tu codigo]
Click \textbf{Run} $\rightarrow$ \textbf{Run Module} (F5) para ejecutar.
\end{notebox}

\section{Paso 5\,---\,Subir tu evaluacion}

Cuando termines:

\begin{enumerate}
  \item Vuelve al script (esta abierto en la barra de tareas).
  \item Click \textbf{Subir Archivos} (boton verde).
  \item Espera 5\,---\,10 segundos mientras hace \texttt{git add}, \texttt{commit} y \texttt{push}.
  \item Aparece un mensaje \textbf{Listo\,---\,Ahora entrega en el AVA} con tu URL.
  \item El URL queda automaticamente \textbf{copiado al portapapeles}.
\end{enumerate}

\section{Paso 6\,---\,Entregar el enlace en el AVA}

\begin{enumerate}
  \item Abre el AVA (Ambiente Virtual de Aprendizaje).
  \item Ve a la \textbf{Evaluacion Parcial} correspondiente (1, 2, 3 o 4 segun el tipo).
  \item En el campo de entrega, pega el enlace con \texttt{Ctrl+V}.
  \item Click \textbf{Enviar entrega}.
\end{enumerate}

\section{Paso 7\,---\,Limpieza (opcional)}

El script te preguntara: \emph{\textquotedblleft\,Ya terminaste la evaluacion\,?\,\textquotedblright}

\begin{itemize}
  \item Si responde \textbf{Si}, elimina la carpeta local del Escritorio. Tu evaluacion queda guardada en GitHub.
  \item Si responde \textbf{No}, la carpeta se conserva por si necesitas seguir editando.
\end{itemize}

% ============================================================
\chapter{Problemas comunes}

\section{\textquotedblleft Windows protegio su PC\textquotedblright}

Primera vez normal. Click \textbf{Mas informacion} $\rightarrow$ \textbf{Ejecutar de todas formas}.

El script desbloquea automaticamente los archivos al iniciar para evitar que esta advertencia se repita.

\section{\textquotedblleft winget no se reconoce\textquotedblright}

Tu Windows no tiene \texttt{winget}. Instala \textbf{App Installer} desde Microsoft Store, o descarga manualmente:

\begin{itemize}
  \item Git: \url{https://git-scm.com/download/win}
  \item GitHub CLI: \url{https://cli.github.com/}
\end{itemize}

\section{\textquotedblleft No tienes una sesion de GitHub activa\textquotedblright}

Click \textbf{Si} cuando el script ofrezca iniciar sesion. Sigue las instrucciones del dialogo con el codigo.

\section{\textquotedblleft No aparece mi tarea en el banner\textquotedblright}

Verifica:

\begin{enumerate}
  \item Tu seccion en el dropdown superior coincide con tu inscripcion real.
  \item El profesor ya activo la tarea en el panel admin (preguntale).
  \item Tienes conexion a internet (el script consulta el backend cada 60 segundos).
\end{enumerate}

\section{\textquotedblleft No aparece mi repositorio al refrescar\textquotedblright}

Verifica:

\begin{enumerate}
  \item Aceptaste el assignment en GitHub Classroom.
  \item Estas logueado con la \textbf{misma cuenta} con la que aceptaste.
  \item Click \textbf{Refrescar} de nuevo.
\end{enumerate}

\section{\textquotedblleft Falla el push\textquotedblright}

Posibles causas:

\begin{itemize}
  \item Conexion a internet caida.
  \item Token expirado: usa el boton \textbf{Cerrar sesion} y vuelve a iniciar sesion.
\end{itemize}

\section{Quiero cambiar de cuenta de GitHub}

Click \textbf{Cerrar sesion} (boton rojo en la esquina superior derecha). Luego \textbf{Iniciar sesion} con la nueva cuenta.

\section{ALERTA roja gigante en mi pantalla}

\begin{warnbox}[NO PANIQUEES]
Si te aparece una ventana roja gigante con el texto \textbf{ALERTA\,---\,POSIBLE TRAMPA ACADEMICA DETECTADA}, significa que el repositorio que intentaste clonar contiene archivos NO permitidos (codigo pre-existente).

Esto puede pasar si:
\begin{itemize}
  \item Aceptaste un repositorio que ya tenia codigo.
  \item Estas intentando clonar un repo de otra persona.
  \item Por error elegiste un repositorio que no es de la evaluacion.
\end{itemize}

\textbf{Solo el profesor puede desbloquear esta ventana} con su clave personal. Avisa al profesor para que se acerque y libere el PC.
\end{warnbox}

% ============================================================
\chapter{Resumen visual del flujo}

\section{Flujo en casa (una vez)}

\begin{lstlisting}
1. Descargar y extraer ZIP
2. Doble-click Subir-Evaluacion.bat
3. Aceptar instalacion winget (Si)
4. Elegir seccion (001D / 002D / 003D)
5. Iniciar sesion en GitHub (codigo XXXX-XXXX)
6. Click banner "Tienes N tareas"
7. Aceptar tarea (se abre Classroom)
8. NO clonar todavia
9. Cerrar el script
\end{lstlisting}

\section{Flujo en el aula (cada evaluacion)}

\begin{lstlisting}
1. Abrir Subir-Evaluacion.bat
2. Modo "Repositorio existente" -> Refrescar
3. Seleccionar tu repo -> Clonar Repo
4. IDLE abre automatico -> programar
5. Subir Archivos
6. URL copiado -> pegar en AVA
\end{lstlisting}

\section{Tabla de tiempo estimado}

\begin{tabular}{l l}
\hline
\textbf{Actividad} & \textbf{Tiempo} \\
\hline
Setup en casa (primera vez)                & 15\,---\,20 min \\
Aceptar nuevo assignment (por evaluacion)  & 1 min \\
Clonar + IDLE en el aula                   & 1 min \\
Subir archivos + AVA                       & 2 min \\
\hline
\end{tabular}

\vspace{1cm}

\begin{successbox}[Recuerda]
Si haces el setup en casa, el dia del examen tienes \textbf{solo 4 minutos de procedimiento tecnico}. El resto del tiempo es para programar tu evaluacion.
\end{successbox}

\end{document}
